%% ============================================================================
%% Fichier  : annexes/Ann_A_glossaire.tex
%% Titre    : Glossaire
%% Version  : 0.1 -- novembre 2026
%% ============================================================================
%% Définitions abrégées des termes marqués \termecle dans les chapitres.
%% ============================================================================

\chapter{Glossaire}
\label{ann:glossaire}

Ce glossaire reprend, sous une forme abrégée, les définitions données dans
les chapitres. Le renvoi indique le chapitre où le terme est défini.

\begin{description}
    \item[\NEK{} (\textit{Not Enough Knowledge})] Cote réservée au contrôle
          dont l'auditeur ne peut pas établir l'état malgré des diligences
          normales (Ch.~\ref{chap:prologue}, \ref{chap:eval-sci}).
    \item[Analyse] Interprétation tirée de constatations au regard d'un
          critère, avec un niveau de confiance déclaré
          (Ch.~\ref{chap:prologue}).
    \item[Assurance limitée] Confiance plus faible, exprimée sous forme
          négative, après des travaux réduits (Ch.~\ref{chap:audit-public}).
    \item[Assurance raisonnable] Confiance élevée mais non absolue, exprimée
          sous forme positive (Ch.~\ref{chap:audit-public}).
    \item[Audit] Examen méthodique, indépendant et documenté, qui compare
          une situation à des critères et conduit à une opinion motivée
          (Ch.~\ref{chap:audit-public}).
    \item[Changement] Toute modification d'un programme, d'un paramètre ou
          d'une base en production (Ch.~\ref{chap:exploitation}).
    \item[Constatation] Fait vérifiable, établi par l'auditeur et appuyé
          sur une preuve conservée (Ch.~\ref{chap:prologue}).
    \item[Contrôle] Dispositif ou action de l'entité qui empêche ou détecte
          une anomalie (Ch.~\ref{chap:audit-public}).
    \item[Contrôle compensatoire] Revue indépendante qui couvre un cumul de
          profils que l'on ne peut pas éviter (Ch.~\ref{chap:sci}).
    \item[Contrôle correctif, détectif, préventif] Contrôle qui corrige,
          repère après coup, ou empêche une anomalie (Ch.~\ref{chap:sci}).
    \item[Contrôle interne] Processus qui fournit une assurance raisonnable
          sur les opérations, l'information et la conformité, selon le
          COSO~2013 (Ch.~\ref{chap:sci}).
    \item[Contrôles d'application] Contrôles intégrés à une application~:
          saisie, traitement, interface, sortie (Ch.~\ref{chap:sci},
          \ref{chap:applications}).
    \item[Contrôles généraux informatiques] Contrôles de l'environnement des
          applications~: accès, changements, exploitation, acquisition
          (Ch.~\ref{chap:sci}).
    \item[Critère d'audit] Exigence de référence, choisie avant le test
          (Ch.~\ref{chap:audit-public}, \ref{chap:normes}).
    \item[Déficience] Insuffisance d'une composante ou d'un principe qui
          réduit la probabilité d'atteindre les objectifs
          (Ch.~\ref{chap:eval-sci}).
    \item[Déficience majeure] Déficience, ou combinaison, qui réduit
          fortement cette probabilité~; elle interdit de conclure à
          l'efficacité (Ch.~\ref{chap:eval-sci}).
    \item[Échantillon] Partie de la population examinée en détail,
          aléatoire ou raisonnée (Ch.~\ref{chap:preuve}).
    \item[Empreinte] Valeur calculée par une fonction de hachage (SHA-256)
          qui change au moindre changement du fichier (Ch.~\ref{chap:preuve}).
    \item[Fraude] Acte intentionnel destiné à obtenir un avantage indu
          (Ch.~\ref{chap:donnees-fraude}).
    \item[Gestion des changements] Contrôles qui garantissent qu'un
          changement est demandé, testé, approuvé, mis en production par un
          tiers et tracé (Ch.~\ref{chap:exploitation}).
    \item[Inspection] Vérification ponctuelle au nom de l'autorité
          hiérarchique (Ch.~\ref{chap:audit-public}).
    \item[Limitation d'étendue] Partie de l'objectif que la mission n'a pas
          pu établir, déclarée au rapport (Ch.~\ref{chap:constater}).
    \item[Matrice des risques et des contrôles] Liste des risques d'un
          circuit et des contrôles attendus en face, établie avant les tests
          (Ch.~\ref{chap:eval-sci}).
    \item[Normalisation] Transformation d'une donnée avant comparaison
          (Ch.~\ref{chap:applications}).
    \item[Norme professionnelle] Règle qui encadre la conduite d'une mission
          (Ch.~\ref{chap:normes}).
    \item[Plan de continuité d'activité] Document qui fixe comment l'entité
          poursuit ses missions pendant une interruption
          (Ch.~\ref{chap:exploitation}).
    \item[Population] Ensemble des opérations sur lesquelles porte la
          conclusion, prise dans le système (Ch.~\ref{chap:preuve}).
    \item[Profil] Ensemble de droits attribué dans une application
          (Ch.~\ref{chap:sci}).
    \item[Propriétaire d'application] Responsable métier qui décide des
          droits dans une application (Ch.~\ref{chap:acces}).
    \item[Questionnaire de contrôle interne] Questions fermées, tirées de la
          matrice, qui recueillent les contrôles déclarés
          (Ch.~\ref{chap:eval-sci}).
    \item[Référentiel] Ensemble structuré de bonnes pratiques où l'on puise
          des critères (Ch.~\ref{chap:normes}).
    \item[Revue des habilitations] Examen périodique et signé des comptes et
          profils par le propriétaire de l'application (Ch.~\ref{chap:acces}).
    \item[Risque inhérent, risque lié au contrôle] Gravité et probabilité
          d'une anomalie sans contrôle~; probabilité que le contrôle ne
          l'empêche ni ne la détecte (Ch.~\ref{chap:preparer}).
    \item[Séparation des tâches] Répartition des étapes d'une opération
          entre des personnes différentes (Ch.~\ref{chap:sci}).
    \item[Système de contrôle interne] Ensemble organisé des dispositifs de
          contrôle interne d'une entité, évalué comme un tout
          (Ch.~\ref{chap:sci}).
    \item[Test de cheminement] Suivi d'une opération réelle du début à la
          fin du circuit (Ch.~\ref{chap:eval-sci}).
    \item[Totaux de contrôle] Nombre de lignes et sommes rapprochés d'une
          source indépendante pour prouver l'exhaustivité
          (Ch.~\ref{chap:preuve}).
    \item[Triangle de la fraude] Pression, occasion, rationalisation
          (Ch.~\ref{chap:donnees-fraude}).
    \item[Vérification] Activité des vérificateurs du CONSUPE, sur pièces et
          sur place (Ch.~\ref{chap:audit-public}).
    \item[Auditologie] Science appliquée de la réduction de l'incertitude
          sur les missions d'audit, selon l'article A0
          (Ch.~\ref{chap:recherche}).
    \item[Divergence de Kullback-Leibler] Mesure, en bits, de l'écart entre
          les distributions de réponses de deux cycles
          (Ch.~\ref{chap:constater}).
    \item[Entropie] Dispersion des réponses \OUI{}, \NON{}, \NEK{}, en bits,
          au plus $\log_2 3$ (Ch.~\ref{chap:constater}).
    \item[Espace d'audit] Ensemble des objets sur lesquels porte une
          mission (Ch.~\ref{chap:preparer}).
    \item[Fiche d'assertion] Objectif, critère, question, risque et
          conséquence d'un contrôle (Ch.~\ref{chap:preparer}).
    \item[Gabarit canonique] Décomposition en huit niveaux, du domaine à la
          trace (Ch.~\ref{chap:preparer}).
    \item[Mission auditologique] Quintuplet $\langle \mathcal{S},
          \mathcal{A}, \mathcal{I}, \mathcal{G}, \mathcal{T} \rangle$ de
          l'article A0 (Ch.~\ref{chap:recherche}).
    \item[Observation] Développement du constat en sept éléments~: norme,
          constatation, risques, contradictoire, analyse, conclusion,
          recommandations (Ch.~\ref{chap:constater}).
    \item[Quadrant diagnostique] Lecture croisée de $\mu$ et de $U$ en quatre
          régions~; la conformité opaque est la plus dangereuse
          (Ch.~\ref{chap:constater}).
    \item[Taux de conformité ($\mu$), taux d'incertitude ($U$)] Parts des
          réponses \OUI{} et \NEK{} (Ch.~\ref{chap:constater}).
\end{description}
